\newpage
\section{Introdução}
\label{sec:introducao}

\subsection{Contextualização}
\label{subsec:contextualizacao}

A transição para a mobilidade elétrica representa um dos maiores desafios de
escala industrial da atualidade, não apenas do ponto de vista energético, mas
sobretudo do ponto de vista de materiais. A demanda global por baterias de
íon-lítio ultrapassou a marca histórica de 1,0~TWh em 2024 e segue em trajetória
de crescimento acelerado, impulsionada tanto pela eletrificação do transporte
quanto pela expansão dos sistemas de armazenamento estacionário
(\textit{Battery Energy Storage Systems} -- BESS) \cite{McKinsey2026, IEA2025}.
Esse crescimento é acompanhado por uma diversificação intensa das químicas de
bateria -- NMC, LFP, LMFP, sódio-íon e estado sólido --, cada uma respondendo a
diferentes equilíbrios entre densidade energética, custo e segurança de
suprimento de materiais críticos.

Nesse cenário, a etapa de fim de vida das baterias ganha relevância crescente.
Estima-se que o volume global de baterias descartadas atinja aproximadamente
1,7~milhão de toneladas até 2030 \cite{melin2022} \pendente{conferir o valor na
fonte; em 2019, a mesma consultoria estimava 1,2~milhão de toneladas}, o que coloca a reciclagem e
o reaproveitamento de materiais críticos -- lítio, níquel, cobalto, manganês,
ferro e fósforo -- como uma peça central da sustentabilidade da própria
transição energética.

\subsection{Problema de Pesquisa}
\label{subsec:problema}

O Brasil ocupa uma posição particular nesse contexto global. Diferentemente da
Europa, cuja trajetória de eletrificação é liderada por veículos puramente
elétricos (BEV) de longo alcance, e da China, onde predominam BEV e híbridos
\textit{plug-in}, a frota brasileira sinaliza uma preferência consolidada por
rotas híbridas -- em especial o Híbrido-Flex (HEV) e o Híbrido \textit{Plug-in}
(PHEV) --, sustentada pela ampla disponibilidade de etanol e pela infraestrutura
de abastecimento já consolidada no país. Essa especificidade produz um perfil
químico próprio: na amostra de emplacamentos analisada neste trabalho
(Seção~\ref{sec:marketshare}), predomina o LFP, seguido do NMC, com uma parcela
expressiva de NiMH associada aos híbridos convencionais -- composição que difere
tanto da europeia, centrada no NMC, quanto da chinesa, centrada no LFP.

Embora a literatura internacional já disponha de estudos de Análise de Fluxo
de Materiais (MFA) para frotas eletrificadas europeias e chinesas, não se
identificaram estudos equivalentes que considerem a composição química
específica da frota brasileira nem os fluxos associados à fabricação de células
de bateria no país. Essa lacuna é agravada pela fragilidade do arcabouço
regulatório nacional: a legislação vigente -- a Política Nacional de Resíduos
Sólidos (Lei 12.305/2010) e a Resolução CONAMA 401/2008
\cite{pnrs2010,conama2008} -- não foi adaptada às especificidades das baterias
de tração de veículos elétricos, e o país ainda não dispõe de metas
quantitativas de reciclagem para essa categoria, ao contrário de China e União
Europeia.

\subsection{Objetivos}
\label{subsec:objetivos}

\subsubsection*{Objetivo Geral}

Caracterizar a frota eletrificada brasileira quanto às químicas de bateria e às
trajetórias de eletrificação até 2040 e quantificar, por meio de uma Análise de
Fluxo de Materiais e Energia (MEFA), os fluxos associados à fabricação de
células LFP -- a química mais frequente na frota --, como subsídio ao
planejamento da cadeia nacional de baterias e de reciclagem.

\subsubsection*{Objetivos Específicos}

\begin{enumerate}
    \item Estabelecer, a partir do cruzamento de dados de emplacamento da ABVE
    com as especificações técnicas do PBEV, a linha de base da participação
    (\textit{market share}) das químicas de bateria (LFP, NMC e NiMH) nos
    veículos eletrificados emplacados no país;

    \item Construir cenários de eletrificação da frota de veículos leves
    (Conservador e Acelerado) para os horizontes de 2030, 2035 e 2040, com base
    no estudo ANFAVEA/BCG e na verificação cruzada com as projeções do PDE 2035,
    da Bright Consulting e do MBC Brasil;

    \item Quantificar, por meio de uma MEFA de fronteira \textit{gate-to-gate},
    os fluxos de materiais, utilidades, emissões e resíduos da fabricação de
    células LFP, por kWh de célula e normalizados para 1~GWh de produção anual;

    \item Discutir as implicações dos resultados para a cadeia nacional de
    reciclagem, com ênfase na sucata de produção e nos gargalos econômicos e
    regulatórios identificados na literatura.
\end{enumerate}

A modelagem dinâmica dos estoques e fluxos de baterias da frota em uso, em fim
de vida e em segunda vida -- prevista no planejamento inicial do projeto -- não
integra o escopo deste relatório e é apresentada como continuidade do trabalho
no Capítulo~\ref{sec:conclusao}.

\subsection{Justificativa e Estrutura do Trabalho}
\label{subsec:justificativa}

A relevância deste trabalho decorre de três fatores complementares: (i) a
dependência brasileira de minerais críticos frequentemente importados, cuja
recuperação doméstica pode mitigar riscos de suprimento; (ii) o passivo
ambiental representado pelo descarte inadequado de baterias em fim de vida; e
(iii) a janela regulatória em aberto -- com projeto de lei em tramitação no
Senado Federal para instituir uma política nacional de circularidade das
baterias veiculares \cite{dialogueearth2026} --, que torna urgente a produção de
evidência técnica que subsidie a formulação dessa política.

O trabalho articula duas escalas de análise. Na escala da frota, caracteriza-se
o mercado brasileiro de veículos eletrificados quanto às químicas de bateria e
às trajetórias de eletrificação; na escala industrial, quantificam-se os fluxos
de materiais e energia de uma fábrica de células LFP, química mais frequente na
amostra analisada. As duas escalas se conectam pela química LFP e pela cadeia de
reciclagem, à qual a fábrica entrega um fluxo de material secundário homogêneo e
disponível desde o início da produção: a sucata de fabricação.

O relatório está estruturado da seguinte forma: o
Capítulo~\ref{sec:metodologia} descreve a metodologia das quatro frentes de
trabalho -- revisão bibliográfica, levantamento de \textit{market share},
construção dos cenários de frota e MEFA da fábrica; o
Capítulo~\ref{sec:revisao} revisa as tecnologias de bateria, a manufatura de
células e as rotas de reciclagem, segunda vida e regulação; o
Capítulo~\ref{sec:resultados} apresenta o \textit{market share}, os cenários de
frota e a transição para o estudo industrial; o Capítulo~\ref{sec:gigafabrica}
apresenta a MEFA da fábrica de células LFP; e o Capítulo~\ref{sec:conclusao}
sintetiza os resultados por objetivo, discute as limitações e aponta a
continuidade do trabalho.
